\section{30.04.2026 - Studium und dann?}

\subsection{Neue Wörter}

\begin{vocabtable}
\deword{der Einstiegsjob}{entry-level job}{noun}
\deword{eher}{rather / more likely}{adverb}
\deword{der Absolvent / die Absolventin}{graduate (male / female)}{noun}
\deword{befristet}{fixed-term / temporary (limited in time)}{adjective}
\deword{der Verantwortungsbereich}{area of responsibility}{noun}
\deword{langfristig}{long-term}{adjective}
\deword{die Frist}{deadline / time limit}{noun}
\deword{sich bieten}{to present itself / to arise (opportunities)}{verb}
\deword{die Weile}{while / short time (rare, mostly in fixed expressions)}{noun}
\deword{der Besitzer}{owner}{noun}
\deword{quatschen}{to chat / talk casually}{verb (colloquial)}
\deword{der Mangel}{lack / shortage / deficiency}{noun}
\deword{wählerisch}{picky / selective}{adjective}
\deword{bisherig}{previous / so far / up to now}{adjective}
\deword{klappen}{to work out / to succeed}{verb}
\deword{der Anteil}{share / proportion / percentage}{noun}
\deword{der Dienst}{service / duty / shift (context-dependent)}{noun}
\deword{sich erholen von + dat}{to recover / to rest / to relax}{verb}
\deword{ehrenamtlich}{voluntary / unpaid}{adjective}
\deword{die Auszeit}{time off / break}{noun}
\deword{der Zweck}{purpose / aim}{noun}
\deword{sich engagieren für + akk}{to get involved / to commit oneself}{verb}
\end{vocabtable}

\subsection{Eigene Sätze}

\begin{itemize}
    \item
    \item 
    \item 
\end{itemize}

\subsection{Schriftlicher Text – Korrektur}


\wrong{Bisherig} \correct{Bisher} meine Erfahrungen sind \wrong{"mixed"} \correct{gemischt}. 
\wrong{Letzte} \correct{Letztes} Jahr hatte ich ein \wrong{Vorstellungsgespraech} \correct{Vorstellungsgespräch} 
und \wrong{seit dann} \correct{seitdem} habe ich eine Arbeit.
\wrong{However,} \correct{Allerdings} habe ich mich auch 
\wrong{fur sehr zu viele Jobs bewoerben} \correct{auf sehr viele Jobs beworben}, 
\wrong{ohne antword} \correct{ohne eine Antwort} 
oder \wrong{sie hat mir} \correct{ich habe} 
der \wrong{klassik} \correct{klassische} 
\enquote{\wrong{leider mussen wir dir lassen lernen dass} 
\correct{Leider müssen wir Ihnen mitteilen, dass …}}


Glassdoor ist auch eine Website, \wrong{wo} \correct{auf der} man sich direkt bewerben kann. 
\wrong{Es ist ausgleich mit SetpStone} \correct{Sie ist vergleichbar mit StepStone}, 
aber \wrong{die einzige Unterschied} \correct{der einzige Unterschied}, 
\wrong{so viel ich weiss} \correct{soweit ich weiß}, ist, 
dass \wrong{mit Glassdoor kann man ihre "preferences einstellen und email notifications bekommen} 
\correct{man bei Glassdoor seine Präferenzen einstellen und E-Mail-Benachrichtigungen erhalten kann}.


\section{30.04.2026 - Warum möchten Sie bei uns arbeiten?}


\subsection{Neue Wörter}

\begin{vocabtable}
\deword{insgesamt}{overall / in total}{Insgesamt fünf Personen.}
\deword{verlaufen}{to proceed / to run; also: to get lost}{Das Meeting verläuft gut.}
\deword{jeweils}{respectively / each}{Die Gruppen arbeiten jeweils separat.}
\deword{die Kompetenz}{skill / competence}{Ich habe viele Kompetenzen.}
\deword{die Kontaktfähigkeit}{interpersonal skills}{Ich arbeite gut im Team.}
\deword{die Führungs- und Verantwortungsbereitschaft}{leadership and responsibility}{Ich übernehme gern Verantwortung.}
\deword{das Durchhaltevermögen}{perseverance}{Ich gebe nicht schnell auf.}
\deword{die Problemlösefähigkeit}{problem-solving ability}{Ich löse Probleme schnell.}
\deword{die Kritikfähigkeit}{ability to accept criticism}{Ich akzeptiere Kritik gut.}
\deword{die Lernbereitschaft}{willingness to learn}{Ich lerne gern Neues.}
\deword{die Vorbereitungsschritte}{preparation steps}{Die Vorbereitungsschritte sind wichtig.}
\deword{erwähnen + akk}{to mention}{Ich erwähne das kurz.}
\deword{der Beitrag}{post / contribution}{Ich schreibe einen Beitrag.}
\deword{üben}{to practice}{Ich übe jeden Tag.}
\deword{begrenzen + akk}{to limit / to restrict}{Ich begrenze meine Zeit.}
\deword{vertiefen + akk}{to deepen / to intensify}{Ich vertiefe mein Wissen.}
\deword{sich fürchten vor + dat}{to fear}{Ich fürchte mich nicht.}
\deword{anstellen + akk}{to hire / to employ}{Die Firma stellt neue Leute an.}
\deword{begründen + akk}{to justify / to explain}{Ich begründe meine Meinung.}
\deword{statt + gen}{instead of}{Ich nehme Wasser statt Cola.}
\deword{bisherig}{previous / so far}{In meinem bisherigen Job habe ich viel gelernt.}
\deword{vorkommen}{to occur / to happen}{ So ein Fehler ist noch nie vorgekommen.}
\deword{begegnen + dat}{to encounter / to come across}{Ich bin dieser Frage schon einmal begegnet.}
\deword{überzeugen + akk}{to convince}{Seine Argumente haben mich überzeugt.}
\deword{die Behauptung}{claim / assertion}{Seine Behauptung war nicht korrekt.}
\deword{die Forderung}{demand / request}{Die Forderungen der Studenten waren klar.}
\deword{die Begründung}{justification / reasoning}{Die Begründung war sehr überzeugend.}
\deword{die Stützung}{support / reinforcement}{Die Daten dienen zur Stützung der Theorie.}
\deword{die Erläuterung}{explanation / clarification}{Die Lehrerin gab eine kurze Erläuterung.}
\deword{derzeit}{currently / at the moment}{Ich arbeite derzeit an meiner Masterarbeit.}

\end{vocabtable}

\subsection{Text – Korrektur}

- Ich finde das Umfrageergebnis nicht überraschend, denn \wrong{meistens Studierenden moecthen im eine grosses Unternehmen arbeiten}
\correct{die meisten Studierenden möchten in einem großen Unternehmen arbeiten}. 
\wrong{Ich bin der Meinung nach das Grund fur diese whal ist wirklich das gehalt}
\correct{Ich bin der Meinung, dass der Grund für diese Wahl vor allem das Gehalt ist}. 
\wrong{Offentlich Internationales Unternehem beiten mehr kompetitiv salaries an}
\correct{Internationale Unternehmen bieten oft \textbf{wettbewerbsfähigere(more competitive)} Gehälter an}.

- \wrong{Warum soll ich dir hire?}
\correct{Warum soll ich dich einstellen?}

- Was ich sehr hilfreich fand\wrong{ ist}\correct{, ist}, dass man \wrong{die Regelmaesig Fragen dass man in einem Vorstellungsgespraech treffen kann}
\correct{regelmäßig Fragen übt, die man in einem Vorstellungsgespräch bekommen kann}. 
Ich hatte \wrong{eine klein} \correct{eine kleine} Idee \wrong{uber} \correct{darüber}, \wrong{was kann man} \correct{was man} in einem Vorstellungsgespräch erwarten kann, 
aber die folgende \wrong{frage} \correct{Frage} \enquote{Welche beruflichen Ziele möchten Sie in \wrong{3} \correct{drei} Jahren erreicht haben?} 
\wrong{habe mir nicht in der Sinn bekommen} \correct{ist mir nicht in den Sinn gekommen}. 
Ich \wrong{bein} \correct{bin} der Meinung, dass man \wrong{muss ihre Dokumente wirklich gut wissen}
\correct{seine Unterlagen wirklich gut kennen muss} und \wrong{fur} \correct{für} alle \wrong{ihre erwaehende Fachickeiten}
\correct{erwähnten Fähigkeiten} sicher sein sollte. 
\wrong{Von meine letze Bewerbungsgespräch} \correct{Aus meinem letzten Bewerbungsgespräch} 
\wrong{hab ich} \correct{habe ich} gelernt, \wrong{mehr zuversichtlich zu sein und mit sichere Worte an alle Fragen antworten}
\correct{zuversichtlicher zu sein und alle Fragen mit sicheren Worten zu beantworten}. 
Leider kann ich keine kulturellen \wrong{Vergleichen} \correct{Vergleiche} \wrong{zwissen} \correct{zwischen} 
\wrong{Deutchland} \correct{Deutschland} und \wrong{Griechenalnd} \correct{Griechenland} \wrong{erwahen} \correct{erwähnen}, 
weil ich kein Vorstellungsgespräch in Griechenland \wrong{hat} \correct{hatte}. 
\wrong{Fur} \correct{Für} neue Studierende und Bewerberinnen \wrong{, waere es ratsam}
\correct{wäre es ratsam}, \wrong{fuer die Job the correct fit zu sein}
\correct{gut zur Stelle zu passen} und \wrong{un dass muss ihre Lebenslauf un Auschreiben ( cover letter) deutlich erklaeren}
\correct{und dies im Lebenslauf und im Anschreiben deutlich zu erklären}.

-\wrong{"Warum moechten Sie hier arbeiten?".}
\correct{\enquote{Warum möchten Sie hier arbeiten?}}
Diese ist eine \wrong{haeufig frage}
\correct{häufige Frage}, die jeder Bewerber erwarten muss.
Zum \wrong{glueck}
\correct{Glück}, \wrong{es hat mir nich passiert diese frage zu entreffen}
\correct{ist mir nicht passiert, dieser Frage zu begegnen}, aber ich bin der Meinung\wrong{ das}
\correct{, dass} es sehr einfach ist\wrong{ auf}
\correct{, auf} diese Frage zu antworten.
Man muss nur\wrong{,}
\correct{} \wrong{das Website des Unternemmen}
\correct{die Website des Unternehmens} lesen und \wrong{dan}
\correct{dann} \wrong{sind Sie bereits}
\correct{ist man bereits vorbereitet}!

-Politische Debatten – Politiker wollen \wrong{die Volk}
\correct{das Volk} über ihre Ideen und \wrong{Plannen}
\correct{Plänen} überzeugen.

-Derzeit arbeite ich bei DLR und bin \wrong{ich verantwortlich frue}
\correct{ich für} die Entwicklung von vier Python-Packages \wrong{}
\correct{verantwortlich}.
